
\subsubsection{Parameter-Efficient Fine-Tuning}

LoRA freezes pre-trained weights W and learns low-rank updates $\Delta W$ = BA where B $\in \mathbb{R}^(d\times r)$, A $\in \mathbb{R}^(r\times d)$, with r typically 8-64. The method demonstrates that task-specific adaptations can be captured in low-dimensional subspaces. LoRA does not explicitly claim that pre-trained weights are low-rank; it shows that updates can be low-rank. Our work measures the property that LoRA does not claim: the rank of W itself.

Subsequent methods including AdaLoRA and DyLoRA explore adaptive rank allocation, indicating that optimal rank varies by task and layer. The existence of rank-search methods suggests low-rank sufficiency is task-dependent rather than a fixed property of weights.

\subsubsection{State Space Models and Hybrid Architectures}

Mamba introduced selective state space models achieving linear-time complexity O(L) with performance matching Transformers at certain scales. Hybrid architectures such as Samba interleave SSM blocks with attention layers, demonstrating complementary strengths when co-trained from initialization.

Post-hoc conversion of pre-trained Transformers to SSMs remains challenging. Our measurement provides one explanation: pre-trained Transformer weights do not exhibit the bounded-state low-rank structure (r\_eff < 256) that certain SSM conversion approaches assume.

\subsubsection{Post-Hoc Compression}

Magnitude pruning, structured pruning, and quantization methods often assume redundancy in pre-trained weights. Knowledge distillation achieves behavioral equivalence without requiring structural similarity. Our finding that projection weights maintain high effective rank (r\_eff ~ 1600) does not preclude these methods but does inform approaches that explicitly assume low-rank weight structure.
